\section*{Introduzione}
\addcontentsline{toc}{section}{Introduzione}

Due popolazioni che coesistono su una rete di traffico possono dare orgine a diversi scenari di interesse. Per riuscire a comprenderli ci serviremo del 
modello presentato nella \autoref{sec:Modello}, ripreso da \cite{MR4911797}. La distinzione nelle popolazioni, denominate
\textit{hat} e \textit{check}, permette di rappresentare due gruppi che possono avere comportamenti simili o del tutto distinti l'uno dall'altro.
Questa differenza viene caratterizzata da due tipi di parametri: i \textit{percorsi} che i veicoli appartenenti ad una delle due popolazioni possono utilizzare per raggiungere la loro destinazione
ed i \textit{costi} delle strade che costituiscono tali percorsi. I costi possono rappresentare, per esempio, il tempo necessario a percorrere la strada, il consumo di benzina oppure l'emissione di inquinanti.

I guidatori scelgono il percorso con un unico obiettivo: minimizzare il costo impiegato nel viaggio sulla rete. Ognuno di essi, dunque,
sceglie il percorso più conveniente per sé stesso, portando la rete ad un \textit{Equilibrio di Nash} \cite{nash1950noncooperative}. In questo stato di equilibrio tutti 
i veicoli in una popolazione pagano lo stesso costo per arrivare a destinazione; nessun guidatore è, inoltre, disposto a cambiare percorso, in quanto vedrebbe
aumentare il proprio costo di percorrenza della rete.

Analizzeremo, dunque, alcune reti che inducono la comparsa di \textit{Paradossi di Braess} \cite{BraessParadox};
queste situazioni paradossali vengono causate dall'aggiunta di strade che aumentano i costi di attraversamento della rete per una o per entrambe le popolazioni.

Nonostante il modello permetta la rappresentazione di diversi tipi di reti, esso non consente l'adeguata descrizione di un caso particolare, discusso nella
\autoref{sec:Critica}.

Viene introdotto uno strumento software a supporto dello studio delle diverse reti; è possibile ottenere grafici rappresentanti i costi di
viaggio sulla rete e gli Equilibri di Nash in funzione di variazioni nel numero di veicoli o nei costi delle singole strade. 